% Copyright (c) 2026 Geovana Neves. Licensed under CC BY 4.0 (https://creativecommons.org/licenses/by/4.0/). See LICENSE-AND-AUTHORSHIP.md.
%
\section{The study's own words}

\begin{frame}[fragile]{The aliases table}
\small An alias is a name YOUR study gives to a set of boundaries, read
everywhere a boundary is cited: a \code{MOVING\_BC\_ALIAS}, a \code{ROTATE}
record, a \code{[groups]} member of the pproc, a \code{families} entry.
\begin{lstlisting}[basicstyle=\ttfamily\scriptsize,numbers=none]
[aliases]
airframe = ["Body", "Base"]
rotors = ["PORT", "STARBOARD"]
\end{lstlisting}
\begin{itemize}\small
  \item A member may be another alias, followed to the end: \code{rotors}
        above names two rotor blocks, and a row can write \code{rotors} to
        reach both.
  \item A member the opened mesh does not carry is \textbf{left out}
        rather than refused. This is what lets one reference file serve a
        wing-body mesh and an isolated-rotor mesh alike.
  \item Declare the set here once; every row, every group and every plot
        then says \code{airframe} instead of listing the boundaries again.
\end{itemize}
\setupsources{\code{cases/\_\_init\_\_.py}, the reference model;
  \code{cases/workflows.py}}
\end{frame}

\begin{frame}[fragile]{Custom coordinate systems: \code{[[frames]]}}
\small One entry per frame, in the order written.
\begin{lstlisting}[basicstyle=\ttfamily\scriptsize,numbers=none]
[[frames]]
name = "NACELLE"          # the name a row's ROTATE axis cites
origin = [4.5, 0.0, 0.0]  # in the geometry's own frame, metres
x_axis = [1.0, 0.0, 0.0]  # optional; default shown
y_axis = [0.0, 1.0, 0.0]  # optional; the third axis is their
                           # right-handed cross product
\end{lstlisting}
\begin{itemize}\small
  \item Created right after the frames the package makes itself
        (\code{MRP}, then \code{<ALIAS>\_SMRP} on a rotor run type), in the
        order written, and before any motion: a rotor whose axis is one of
        these turns about a frame that already exists.
  \item A row's rotation or translation cites a frame axis as
        \code{<frame>-X}, \code{<frame>-Y} or \code{<frame>-Z}.
  \item Refused at plan time: a name the package already creates, a name
        defined twice, or an origin that is not three numbers.
\end{itemize}
\setupsources{\code{cases/\_\_init\_\_.py}, the reference model;
  \code{cases/workflows.py}}
\end{frame}
